

\chapter{Ergodic theorems}

\section{Von Neumann's $L_2$ ergodic theorem}
Our first ergodic theorem is von Neumann's ergodic
theorem, which in fact is a result in operator theory
that has a nice interpretation for our problem of the
convergence of ergodic averages in a measure
preserving system.

\begin{thm}[Von Neumann's ergodic theorem]
\label{thm-vonneumann}
Let $H$ be a Hilbert space, and let $U$ be a unitary
operator on $H$. Let $P$ be the \emph{orthogonal} projection
operator onto the subspace $\operatorname{Ker}(U-I)$
(the subspace of $H$ consisting of $U$-invariant
vectors). For any vector $v\in H$ we have

\begin{equation*} \tag{17}\label{eq:vonneumann}
\frac{1}{n}\sum_{k=0}^{n-1} U^k v \to Pv \ \textrm{ as }n\to\infty.
\end{equation*}

\noindent (Equivalently, the sequence of operators
$\frac{1}{n}\sum_{k=0}^{n-1} U^k$ converges to $P$ in
the strong operator topology.)
\end{thm}

\begin{proof}
Define two subspaces

\begin{align*}
V&=\operatorname{Ker}(U-I) = \{ v\in H\,:\, Uv=v \}, \\
V'&=\operatorname{Range}(U-I) = \{ Uw-w \,:\, w\in H \}.
\end{align*}

\noindent Note that (\ref{eq:vonneumann}) holds trivially for
$v\in V$. For a different reason, we also show that
it holds for $v\in V'$: if $v=Uw-w$ then we have

\[
\frac{1}{n}\sum_{k=0}^{n-1} U^k v  = \frac{1}{n}\left( U^n w - w \right) \to 0\ \textrm{ as }n\to\infty.
\]

\noindent On the other hand, one can verify that $v\in V^\perp$,
and therefore $Pv=0$, by observing that if $z\in V$
then

\[
\langle w, z \rangle = \langle Uw,Uz \rangle =
\langle Uw,z\rangle,
\]

\noindent hence $\langle Uw-w,z\rangle = 0$.

Combining the above observations we see that
(\ref{eq:vonneumann}) holds for $v\in V+V'$. Next, we
claim that it also holds for $v\in \overline{V+V'}$,
the norm closure of $V+V'$. Indeed, if
$v\in \overline{V+V'}$ then for an arbitrary
$\epsilon>0$ we can take $w\in V+V'$ such that
$\lVert v-w\rVert < \epsilon$ and conclude that

\begin{align*}
\left\lVert \left(\frac{1}{n}\sum_{k=0}^{n-1} U^k-P\right)v \right\rVert
& \le
\left\lVert \left(\frac{1}{n}\sum_{k=0}^{n-1} U^k-P\right)w \right\rVert
+
\left\lVert \left(\frac{1}{n}\sum_{k=0}^{n-1} U^k-P\right)(v-w) \right\rVert \\
& \le
\left\lVert \left(\frac{1}{n}\sum_{k=0}^{n-1} U^k-P\right)w \right\rVert
+  \epsilon.
\end{align*}

\noindent This implies that
$\limsup_{n\to\infty} \left\lVert
\left(\frac{1}{n}\sum_{k=0}^{n-1} U^k-P\right)v
\right\rVert < \epsilon$,
and since $\epsilon$ was an arbitrary positive number
we get (\ref{eq:vonneumann}).

Finally, we claim that $H=\overline{V+V'}$. Since
$\overline{V+V'}$ is a \emph{closed} subspace of $H$, we
have

\[
\overline{V+V'} = \left( \left(\overline{V+V'}\right)^\perp \right)^\perp,
\]

\noindent (in general, the orthogonal complement of the
orthogonal complement of a subspace $W$ of a Hilbert
space is equal to $\overline{W}$). So, it suffices to
show that $\left(\overline{V+V'}\right)^\perp=\{0\}$,
i.e., that the only vector orthogonal to all of
$\overline{V+V'}$ is the zero vector. Assume $w$ is
such a vector. Then $w \perp Uw-w$. But note that we
have the identity

\begin{align*}
\lVert Uw-w \rVert^2 &= \langle Uw-w, Uw-w\rangle = \lVert Uw \rVert^2 + \lVert w \rVert^2 - 2\operatorname{Re} \langle Uw,w\rangle  \\ &=
2\lVert w \rVert^2 - 2\operatorname{Re} \langle Uw,w\rangle
= - 2\operatorname{Re} \langle Uw-w,w\rangle
\end{align*}

\noindent which means that $Uw-w = 0$, i.e., $w\in V$. Since
$w\in \left(\overline{V+V'}\right)^\perp$ we get that
$w$ is orthogonal to itself and therefore $w=0$, as
claimed.
\end{proof}

\noindent Let $(\Omega, \mathcal{F}, \mathrm{P}, T)$ be a
measure preserving system. We associate with $T$ an
operator $U_T$ on the Hilbert space $L_2(\Omega)$,
defined  by

\[
U_T(f) = f\circ T.
\]

\noindent The fact that $T$ is measure preserving implies that
$U_T$ is unitary:

\[
\langle U_T f, U_T g \rangle = \mathbf{E}( (U_T f) \overline{(U_T g)} )
= \mathbf{E}( (f\circ T) \overline{(g\circ T)} )
= \mathbf{E}( (f\overline{g})\circ T ) = \mathbf{E}( f\overline{g} ) = \langle f,g\rangle.
\]

\noindent Note also that the subspace $\operatorname{Ker}(U-I)$
consists exactly of the invariant (square-integrable)
random variables, or equivalently those random
variables which are measurable with respect to the
$\sigma$-algebra $\mathcal{I}$ of invariant events.
Recalling the discussion of conditional expectations
in Chapter 2, we also see that the orthogonal
projection operator $P$ is exactly the conditional
expectation operator
$\mathbf{E}\left(\cdot\,|\,\mathcal{I}\right)$ with
respect to the $\sigma$-algebra of invariant events!
Thus, Theorem \ref{thm-vonneumann} applied to this setting
gives the following result.

\begin{thm}[The $L_2$ ergodic theorem]
Let $(\Omega, \mathcal{F}, \mathrm{P}, T)$ be a
measure preserving system. For any random variable
$X\in L_2(\Omega, \mathcal{F}, \mathrm{P})$, we have

\[
\frac{1}{n} \sum_{k=0}^{n-1} X\circ T^{k} \to \mathbf{E}\left(X\,|\,\mathcal{I}\right)\ \ \textrm{in $L_2$ as }n\to\infty.
\]

\noindent In particular, if the system is ergodic then

\[
\frac{1}{n} \sum_{k=0}^{n-1} X\circ T^{k} \to \mathbf{E}(X)\ \ \textrm{in $L_2$ as }n\to\infty.
\]

\end{thm}

\section{Birkhoff's pointwise ergodic theorem}
We will now prove \emph{the} fundamental result of ergodic
theory, known alternately as
$\textbf{Birkhoff's pointwise ergodic theorem}$;
$\textbf{Birkhoff's ergodic theorem}$; the
$\textbf{pointwise ergodic theorem}$; or just the
$\textbf{ergodic theorem}$\footnote{Incidentally, I've always found it strange that
ergodic theory --- unlike other areas of math ---
seems to be the only theory named after an adjective
(as opposed to a noun, as in \emph{the theory of numbers},
or as in the non-existent name \emph{the theory of
ergodicity}, which would perhaps have been a better
name for ergodic theory). Similarly, the ergodic
theorem is, as far as I know, the only theorem in math
to be named after an adjective. (And what does the
name mean, anyway? That the theorem has no nontrivial
invariant sets...?) If you think of any
counterexamples to this observation, please let me
know!}.

\begin{thm}[Birkhoff's ergodic theorem]
Let $(\Omega, \mathcal{F}, \mathrm{P}, T)$ be a
measure preserving system. Let $\mathcal{I}$ denote
as usual the $\sigma$-algebra of $T$-invariant sets.
For any random variable
$X \in L_1(\Omega, \mathcal{F}, \mathrm{P})$, we have

\begin{equation*} \tag{18}\label{eq:birkhoffergodic1}
\frac{1}{n} \sum_{k=0}^{n-1} X\circ T^k \xrightarrow[]{\textrm{\,a.s.\,}} \mathbf{E}\left(X\,|\,\mathcal{I}\right)\ \textrm{ as }n\to\infty.
\end{equation*}

\noindent When the system is ergodic, we have

\begin{equation*} \tag{19}\label{eq:birkhoffergodic2}
\frac{1}{n} \sum_{k=0}^{n-1} X\circ T^k \xrightarrow[]{\textrm{\,a.s.\,}} \mathbf{E}(X)\ \textrm{ as }n\to\infty.
\end{equation*}

\end{thm}

\noindent For the proof, we start by proving a lemma, known as
the $\textbf{maximal ergodic inequality}$.

\begin{lem}
\label{lem-ergodic-maximal}
With the same notation as above, denote also $S_0=0$,
$S_n  = \sum_{k=0}^{n-1} X\circ T^k$, and let
$M_n=\max\{ S_k \,:\,0\le k\le n\}$. For each $n\ge1$
we have

\[
\mathbf{E}(X \mathbf{1}_{\{M_n > 0 \}}) \ge 0.
\]

\end{lem}

\begin{proof}
For each $0\le k\le n$ we have

\[
S_{k+1} = X + S_k\circ T \le X + M_n\circ T,
\]

\noindent or equivalently $X \ge S_{k+1} - M_n\circ T$. Since
this is true for each $0\le k\le n$, we get that

\[
X \ge \max(S_1,\ldots,S_n)-M_n\circ T,
\]

\noindent and therefore, noting that on the event $\{M_n>0\}$,
we have $M_n=\max(S_1,\ldots,S_n)$, we get that

\begin{align*}
\mathbf{E}(X \mathbf{1}_{\{M_n > 0 \}})
&\ge
\mathbf{E}\left[(\max(S_1,\ldots,S_n) - M_n\circ T) \mathbf{1}_{\{M_n > 0 \}}\right] \\
&=
\mathbf{E}\left[(M_n - M_n\circ T) \mathbf{1}_{\{M_n > 0 \}}\right] \\
&=
\mathbf{E}\left[(M_n - M_n\circ T)\right]
-
\mathbf{E}\left[(M_n - M_n\circ T) \mathbf{1}_{\{M_n > 0 \}^c}\right] \\
&= 0 -
\mathbf{E}\left[(M_n - M_n\circ T) \mathbf{1}_{\{M_n = 0 \}}\right] \\
&=
\mathbf{E}\left[(M_n\circ T) \mathbf{1}_{\{M_n = 0 \}}\right] \ge 0.
\end{align*}

\end{proof}

\begin{proof}[Proof of the ergodic theorem]
$\mathbf{E}\left(X\,|\,\mathcal{I}\right)$ is an
invariant random variable, so by replacing $X$ with
$X-\mathbf{E}\left(X\,|\,\mathcal{I}\right)$, we can
assume without loss of generality that
$\mathbf{E}\left(X\,|\,\mathcal{I}\right)=0$; in this
case, we need to prove that $S_n/n\to 0$ almost
surely (where $S_n=\sum_{k=0}^{n-1} X\circ T^k$ as in
the lemma above). Denote
$\overline{X} = \limsup_{n\to\infty} S_n/n$.
$\overline{X}$ is an invariant random variable,
taking values in $\R\cup\{\pm\infty\}$. Fix
$\epsilon>0$, and consider the invariant event
$A= \{ \overline{X} > \epsilon \}$. We claim that
$\mathrm{P}(A)=0$. Once we prove this, since
$\epsilon$ is arbitrary it will follow that
$\overline{X}\le 0$ almost surely. By applying the
same result to $-X$ instead of $X$ the reverse
inequality that almost surely $\liminf S_n/n\ge 0$
will also follow, and the theorem will be proved.

To prove the claim, define a new random variable
$X^* = (X-\epsilon)\mathbf{1}_{A}$. Applying Lemma
\ref{lem-ergodic-maximal} to $X^*$ we get that

\[
\mathbf{E}\left[X^* \mathbf{1}_{\{ M_n^*>0\}}\right] \ge 0,
\]

\noindent where $M_n^* = \max(0,S_1^*,\ldots,S_n^*)$ and

\[
S_k^* = \sum_{j=0}^{k-1} X^*\circ T^j = \sum_{j=0}^{k-1} \left((X-\epsilon)\circ T^{k-1}\right) \mathbf{1}_A = (S_k-k\epsilon)\mathbf{1}_A
\]

\noindent (since $A$ is an invariant event). Note that the
events $\{M_n^* > 0\}$ are increasing, so
$X^* \mathbf{1}_{\{ M_n^*>0\}} \to X^* \mathbf{1}_B$
almost surely as $n\to\infty$, where the event $B$ is
defined by

\[
B=\bigcup_{n=1}^\infty \{M_n^* > 0 \} =
\left\{ \sup_{n\ge 1} S_n^* > 0 \right\}
= \left\{ \sup_{n\ge 1} S_n^*/n > 0 \right\}.
\]

\noindent Furthermore, the convergence is dominated, since
$\mathbf{E}|X^*| \le \mathbf{E}|X|+\epsilon<\infty$,
so the dominated convergence theorem implies that

\[
\mathbf{E}(X^* \mathbf{1}_B) \ge 0.
\]

\noindent Finally, observe that $A\subset B$, because

\begin{align*}
A&=\{\limsup_{n\to\infty} S_n/n>\epsilon\} \subseteq A \cap \{ S_n>n\epsilon \textrm{ for some }n\ge 1\} \\
&= \bigcup_{n=1}^\infty \{ (S_n-n\epsilon)\mathbf{1}_A > 0 \} = \left\{ \sup_{n\ge 1} S_n^*>0 \right\} = B.
\end{align*}

\noindent So we have shown that

\begin{align*}
0 &\le \mathbf{E}(X^*\mathbf{1}_B) = \mathbf{E}((X-\epsilon)\mathbf{1}_A \mathbf{1}_B) =
\mathbf{E}((X-\epsilon)\mathbf{1}_{A\cap B} ) = \mathbf{E}((X-\epsilon)\mathbf{1}_A ) \\&=
\mathbf{E}(X\mathbf{1}_A) - \epsilon \mathrm{P}(A)
= \mathbf{E}(\mathbf{E}\left(X\,|\,\mathcal{I}\right)\mathbf{1}_A) - \epsilon \mathrm{P}(A) = -\epsilon \mathrm{P}(A),
\end{align*}

\noindent which proves our claim that $\mathrm{P}(A)=0$.
\end{proof}

\section{The $L_1$ ergodic theorem}
A trivial addendum to the previous proof shows that
we also get $L_1$ convergence of the ergodic
averages.

\begin{thm}[$L_1$ ergodic theorem]
The convergence in (\ref{eq:birkhoffergodic1}),
(\ref{eq:birkhoffergodic2}) is also in $L_1$.
\end{thm}

\begin{proof}
Fix $M>0$, and write $X=Y_M + Z_M$ where
$Y_M = X \mathbf{1}_{\{|X|\le M\}}$ and
$Z_M = X-Y_M = X \mathbf{1}_{\{|X|>M\}}$. The
pointwise ergodic theorem implies that

\[
\frac{1}{n} \sum_{k=0}^n Y_M \circ T^k \to \mathbf{E}\left(Y_M\,|\,\mathcal{I}\right)
 \ \ \textrm{almost surely as }n\to\infty,
\]

\noindent and since $|Y_M|\le M$ the bounded convergence
theorem implies also convergence in $L_1$, i.e.,

\begin{equation*} \tag{20}\label{eq:ergodic-trunc-l1}
\left|\frac{1}{n} \sum_{k=0}^n Y_M \circ T^k - \mathbf{E}\left(Y_M\,|\,\mathcal{I}\right)\right| \to 0 \ \ \textrm{ as }n\to\infty.
\end{equation*}

\noindent Next, for $Z_M$ we have the trivial estimates

\begin{align*}
\mathbf{E} \left| \frac{1}{n} \sum_{k=0}^{n-1} Z_M \circ T^k \right|
& \le \sum_{k=0}^{n-1} \mathbf{E}|Z_M \circ T^k| = \mathbf{E}|Z_M|, \\
\mathbf{E} \left| \mathbf{E}\left(Z_M\,|\,\mathcal{I}\right) \right| \le \mathbf{E} \mathbf{E}\left(|Z_M|\,|\,\mathcal{I}\right) = \mathbf{E}|Z_M|,
\end{align*}

\noindent so, combining this with (\ref{eq:ergodic-trunc-l1}), this
shows that almost surely

\[
\limsup_{n\to\infty} \mathbf{E}\left| \frac{1}{n} \sum_{k=0}^{n-1} X\circ T^k - \mathbf{E}\left(X\,|\,\mathcal{I}\right)\right|
\le 2 \mathbf{E}|Z_M|.
\]

\noindent Letting $M\to\infty$ finishes the proof, since
$\limsup_{M\to\infty} \mathbf{E}|Z_M| = 0$ by the
dominated convergence theorem.
\end{proof}

\begin{xexercise}
Prove that if $X$ is in $L_p$ for some $p>1$ then the
convergence in (\ref{eq:birkhoffergodic1})  is also in the
$L_p$ norm.
\end{xexercise}

\section{Consequences of the ergodic theorem}
In probability theory and many related fields, the
ergodic theorem is an essential tool that is used
frequently in concrete situations. Here are some of
its consequences with regards to some of the examples
we discussed before.

\begin{enumerate}[label=\arabic*.]
\item $\textbf{The strong law of large numbers.}$ If
$X_1,X_2,\ldots$ are i.i.d. with
$\mathbf{E}|X_1|<\infty$, then if we think of the
variables as being defined on the canonical product
space $\R^\N$ (i.e., $X_n=\pi_n(\omega)$ is the $n$th
coordinate function), then we have
$X_n = X_1\circ S^{n-1}$, where $S:\R^\N\to\R^\N$ is
the shift map. Thus, the ergodic average
$\frac{1}{n} \sum_{k=0}^{n-1} X_1\circ S^k$ is the
same as the familiar empirical average
$\frac{1}{n} S_n = \frac{1}{n} \sum_{k=1}^n X_k$ for
an i.i.d., sum, and Birkhoff's ergodic theorem implies
the strong law of large numbers. (In fact, one can
think of the ergodic theorem as a powerful and
far-reaching generalization of the SLLN).

\item \textbf{Equidistribution of the fractional part of
$n\alpha$.} A classical question in number theory
concerns the statistical properties of the fractional
part of the integer multiples of a number $\alpha$,
i.e., the sequence $\{n \alpha \}$ (sometimes written
as $n\alpha \textrm{ mod }1$), where
$\{z\}=z-\lfloor z\rfloor$ denotes the fractional
part of a real number $z$. If $\alpha$ is a rational
number, it is easy to see that this sequence is
periodic, and its range is the finite set of numbers
$\left\{ \frac{k}{q}\,:\, k=0,1,\ldots,q-1 \right\}$
(where $q$ is the denominator in the representation
of $\alpha$ as a reduced fraction $p/q$), so the
question is trivial. In the case of irrational
$\alpha$ something nice (though not too surprising,
in hindsight) happens:

\begin{thm}[Equidistribution theorem]
\label{thm-equidistribution}
If $\alpha\in \R\setminus\mathbb{Q}$ then the
sequence $( \{ n\alpha \} )_{n=1}^\infty$ is
\textbf{equidistributed in $[0,1]$}\footnote{This is the terminology used in number theory --- see
for example Section ? in the book \emph{An Introduction to
the Theory of Numbers}, by Hardy and Wright, and the
Wikipedia article
\href{http://en.wikipedia.org/wiki/Equidistributed_sequence}{http://en.wikipedia.org/wiki/Equidistributed\_sequence}.
Note that in probability theory the word
\emph{equidistributed} means equal in distribution rather
than uniformly distributed, so one should take care
when using this term for the number theoretic meaning
when talking to a probabilist.}.
More precisely, for any $0<a<b<1$ we have

\[
\frac{1}{n} \# \Big\{ 1\le k\le n\,:\, \{n\alpha\}\in (a,b) \Big\} \to b-a\ \textrm{ as }n\to\infty.
\]

\end{thm}

\noindent To prove this, note that $\{ n\alpha \}$ is simply
$R_\alpha(0)$, where $R_\alpha$ is the circle
rotation map discussed in previous sections. Since we
proved that $R_\alpha$ is ergodic when $\alpha$ is
irrational, the ergodic theorem implies that for
almost every $x\in[0,1]$

\[
\frac{1}{n} \# \Big\{ 1\le k\le n\,:\, \{x+n\alpha\}\in (a,b) \Big\}
= \frac{1}{n} \sum_{k=0}^{n-1} \left(\mathbf{1}_{(a,b)}\circ R_\alpha^k\right) (x)
\to \int_0^1 \mathbf{1}_{(a,b)}(u)\,du = b-a
\]

\noindent as $n\to\infty$. This would appear to be a weaker
result, since it doesn't guarantee that the
convergence occurs for the specific initial point
$x=0$. However, in the particular example of the
irrational circle rotation map (and the particular
observable of the form $\mathbf{1}_{(a,b)}$) a
slightly unusual thing happens, which is that the
ergodic theorem turns out to be true not just for
almost every initial point $x$ but for \emph{all}$x$; in
fact, it is easy to see that convergence for one value
of $x$ is equivalent to convergence for any other
value of $x$ (and in particular $x=0$). This is left
to the reader as an exercise.

$\textbf{Note.}$ Theorem \ref{thm-equidistribution} was
proved in 1909 and 1910 independently by Weyl,
Sierpinski and Bohl. In 1916 Weyl showed that the
sequence $\{ n^2 \alpha \}$ is equidistributed, and
more generally that $\{ p(n) \}$ is equidistributed
if $p(x)$ is a polynomial with at least one
irrational coefficient. Vinogradov proved in 1935 that
if $\alpha$ is irrational then the sequence
$\{ p_n \alpha \}$ is equidistributed, where $p_n$ is
the $n$th prime number. Jean Bourgain (winner of a
1994 Fields Medal) proved similar statements in the
more general setting of the pointwise ergodic theorem
(i.e., the ergodic averages of the form
$\frac{1}{n}\sum_{k=1}^n X\circ T^{k^2}$ and
$\frac{1}{n}\sum_{k=1}^n X\circ T^{p_k}$ in a measure
preserving system converge almost surely, under mild
integrability conditions).

\item $\textbf{Benford's law.}$ A beautiful variant of the
circle rotation example above involves multiplication
instead of addition (but one then has the luxury of
multiplying by nice numbers such as rational numbers
or integers, instead of adding irrational numbers).
Consider for example the distribution of the first
digit in the decimal expansion of the sequence of
powers of $2$, $(2^n)_{n=1}^\infty$. Should we expect
all digits to appear equally frequently? No, a quick
empirical test shows that small digits appear with
higher frequency than large digits. To see why, note
that this is related to the dynamical system
$T:x\mapsto 2x \textrm{ mod } (10^k)_{k=1}^\infty$ on
the interval $[1,10)$ (i.e., multiplication by $2$ in
the quotient group of all positive numbers with the
multiplication operator quotiented by the cyclic group
generated by the number $10$). For example, starting
from $1$ and iterating the map we get the sequence

\[
1 \mapsto 2 \mapsto 4 \mapsto 8 \mapsto 1.6 \mapsto 3.2 \mapsto 6.4 \mapsto 1.28 \mapsto \ldots
\]

\noindent It is easy to check that this map has the invariant
measure

\[
d\mu(x) = \frac{1}{\log 10} \frac{dx}{x} \qquad (0<x<1)
\]

\noindent In fact, this is a thinly disguised version of the
circle rotation map $R_\alpha$ with
$\alpha=\log_{10} 2$; the two maps are conjugate by
the mapping $C(x) = \log_{10} x$ (i.e., $C$ maps
$[1,10)$ bijectively to $[0,1)$ and the relation
$T = C^{-1} \circ R_\alpha\circ C$ holds), and
furthermore the measure $\mu$ defined above is the
pull-back of Lebesgue measure on $[0,1)$ with respect
to the conjugation map $C$, which is why an
experienced ergodicist will know immediately that
$\mu$ is an invariant measure for
$T$.\footnote{We are skirting an important concept in ergodic
theory here --- in fact, the map $C$ is an example of
an $\textbf{isomorphism}$ between two measure
preserving systems. Isomorphisms play a central role
in ergodic theory, and there's a lot more to say about
them, but we will not go further into the subject due
to lack of time.}

With this setup, we can answer the question posed at
the beginning of the example. For each $1\le d\le 9$,
the fraction of the first $n$ powers of $2$ whose
decimal expansions start with a given digit $d$ is
given by

\begin{align*}
\frac{1}{n} \#\{ 0\le k\le n-1\,&:\, T^k(1) \in [d,d+1) \} =
\frac{1}{n} \sum_{k=0}^{n-1} \left(\mathbf{1}_{[d,d+1)}\circ T^k\right)(1) \\ &\to \frac{1}{\log 10}\int_d^{d+1}\frac{dx}{x}
= \log_{10}\frac{d+1}{d},
\end{align*}

\noindent where the convergence follows by the equidistribution
theorem (Theorem \ref{thm-equidistribution}), the above
comments and the exercise below. This probability
distribution on the numbers $1,\ldots,9$ is known as
$\textbf{Benford's law}$. Note that the most common
digit $1$ appears more than $30\%$ of the time, and
the least frequent digit $9$ only appears only $4.6\%$
of the time.

\begin{center}
\fitwidth{$\displaystyle
\begin{array}{c|c|l}
d & \mathbb{P}_\textrm{Benford}(d) = \log_{10}\tfrac{d+1}{d} & \text{Graphical illustration} \\
\hline
1 & 30.1\% & \rule{30.1mm}{8pt} \\
2 & 17.6\% & \rule{17.6mm}{8pt} \\
3 & 12.5\% & \rule{12.5mm}{8pt} \\
4 & 9.7\% & \rule{9.7mm}{8pt} \\
5 & 7.9\% & \rule{7.9mm}{8pt} \\
6 & 6.7\% & \rule{6.7mm}{8pt} \\
7 & 5.8\% & \rule{5.8mm}{8pt} \\
8 & 5.1\% & \rule{5.1mm}{8pt} \\
9 & 4.6\% & \rule{4.6mm}{8pt}
\end{array}
$}
\end{center}

\noindent \emph{Table: The (approximate) digit frequencies in
Benford's law.}

\begin{xexercise}
Let $n<m$ be positive integers. Prove that if $m$ is
not an integer power of $n$ then $\log_m n$ is an
irrational number.
\end{xexercise}

\item $\textbf{Continued fractions.}$ The fact that the
continued fraction map on $(0,1)$ (together with the
Gauss invariant measure) is ergodic has important
consequences regarding the distribution of quotients
in the continued fraction expansion of a number chosen
uniformly at random in $(0,1)$. In contrast to the
much more trivial case of the digits in a decimal (or
base-$b$) expansion, which are simply i.i.d.\ random
numbers chosen from $0,\ldots,9$, the asymptotic
distribution of successive continued fraction
quotients is that they are identically distributed,
but not quite independent. To see this, note that the
marginal distribution of a single quotient can be
computed using ergodic averages, as follows. For each
$q\ge 1$, the set of numbers $x$ whose first quotient
$N(x)=\lfloor 1/x \rfloor$ is equal to $q$ is exactly
the interval $\left(\frac{1}{q+1},\frac{1}{q}\right]$.
Thus, for a given number $x$ we can recover the
proportion of the first $n$ quotients equal to $q$ as

\begin{align*}
\frac{1}{n} \sum_{k=0}^{n-1} \left(\mathbf{1}_{(1/(q+1),1/q]}\circ G^k\right)(x),
\end{align*}

\noindent which by the ergodic theorem converges to

\begin{equation*} \tag{21}\label{eq:gauss-quotient}
\frac{1}{\log 2}\int_{1/(q+1)}^{1/q} \frac{dx}{1+x} = \log_2 \left( \frac{(q+1)^2}{q(q+2)}\right)
\end{equation*}

\noindent for a set of $x$'s that has measure $1$ with respect
to Gauss measure $\gamma$ (and hence, also almost
surely with respect to Lebesgue measure, since
$\gamma$ and Lebesgue are mutually absolutely
continuous with respect to each other). Thus, the
formula on the right-hand side of
(\ref{eq:gauss-quotient}) (which is oddly reminiscent of
Benford's law, though they are not related) represents
the limiting distribution of the first quotient of a
random number. For example, the frequency of
occurrence of the quotient $1$ is
$\log_2(4/3) \approx 41.5\%$ --- more than $40\%$ of
the quotient are equal to $1$\,! Note that this is an
asymptotic result that pertains to the statistics of
many quotients of a given number $x$, and not to the
\emph{first quotient of $x$}: if $x$ is chosen uniformly
in $[0,1]$, because Lebesgue measure is not invariant
under the Gauss map $G$, the first quotient of $x$
has a different distribution (clearly, the probability
that the first quotient is $q$ is exactly
$1/q-1/(q+1)$, the length of the interval
$(1/(q+1),1/q]$).

\begin{xexercise}
Compute the asymptotic probability that a pair of
successive quotients of a randomly chosen $x$ in
$[0,1]$ is equal to $(1,1)$ and compare this to the
square of the frequency of $1$'s, to see why
successive quotients are not independent of each
other. Are two successive $1$'s positively or
negatively correlated?
\end{xexercise}

\noindent What other quantities of interest can one compute for
the continued fraction expansion of random numbers?
One can try computing the expected value of a
quotient, but that turns out not to be very
interesting --- the average
$\frac{1}{\log 2} \int_0^1 N(x) \frac{dx}{1+x}$ is
infinite. The Russian probabilist Khinchin (known for
his Law of Iterated Logarithm, a beautiful result on
random walks and Brownian motion) derived an
interesting limiting law for the \emph{geometric} average
of the quotients. He proved that for almost every
$x\in[0,1]$, the geometric average
$(q_1\ldots q_n)^{1/n}$ of the first $n$ quotients of
$x$ converges to the constant

\[
K = \prod_{k=1}^\infty \left( 1+\frac{1}{k(k+2)} \right)^{\log_2 k} \approx 2.68545
\]

\noindent (known as $\textbf{Khinchin's constant}$).

\begin{xexercise}
Prove Khinchin's result.
\end{xexercise}

\noindent We mention one additional and very beautiful limiting
result on continued fraction expansions. If
$x\in (0,1)$ has an infinite continued fraction
expansion

\[
x = \frac{1}{n_1+\frac{1}{n_2+\frac{1}{n_3+\frac{1}{n_4+\frac{1}{\ldots}}}}},
\]

\noindent where $n_1,n_2,\ldots$ are the quotients in the
expansion, it is interesting to consider the truncated
expansion

\[
\frac{P_k}{Q_k} = \frac{1}{n_1+\frac{1}{n_2+\frac{1}{n_3+\frac{1}{\ddots+\frac{1}{n_k}}}}},
\]

\noindent which are rational numbers that become better and
better approximations to $x$. In fact, one reason why
continued fraction expansions are so important in
number theory is that it can be shown that the \emph{best}
rational approximation to $x$ with denominator
bounded by some bound $N$ will always be the last
truncated continued fraction $P_k/Q_k$ for which
$Q_k\le N$, and furthermore, the inequalities

\begin{equation*} \tag{22}\label{eq:cont-frac-approx}
\frac{1}{Q_k (Q_k+Q_{k+1})} \le \left|x - \frac{P_k}{Q_k} \right| \le \frac{1}{Q_k Q_{k+1}}
\end{equation*}

\noindent hold. How fast should we expect this sequence of
rational approximations to converge? The answer is
given in the following theorem. For the proof (which
is surprisingly not difficult), see Section~1.4 in the
book \emph{Ergodic Theory and Information} by P.\
Billingsley.

\begin{thm}
For almost every $x\in (0,1)$ we have

\begin{align*}
\lim_{k\to\infty} \frac{1}{k}\log Q_k &= \tag{23}\label{eq:cont-frac-entropy1} \frac{\pi^2}{12\log 2},
\\
\lim_{k\to\infty} \frac{1}{k}\log \left| x-\frac{P_k}{Q_k} \right| &= \tag{24}\label{eq:cont-frac-entropy2} \frac{\pi^2}{6\log 2},
\\
\lim_{k\to\infty} \frac{1}{k}\log \operatorname{Leb}(\Delta_k(x)) &= \tag{25}\label{eq:cont-frac-entropy3} \frac{\pi^2}{6\log 2},
\end{align*}

\noindent where
$\Delta_k(x) = \{ y\in (0,1) : n_j(y)=n_j(x) \textrm{
for }1\le j\le k\}$
(this interval is sometimes called the \textbf{$k$th
fundamental interval of $x$}), and
$\operatorname{Leb}(\cdot)$ denotes Lebesgue measure
(it is easy to see that the same statement is true if
Gauss measure is used instead).
\end{thm}

\noindent Note that (\ref{eq:cont-frac-entropy2}) and
(\ref{eq:cont-frac-entropy3}) follow easily by combining
(\ref{eq:cont-frac-entropy1}) with
(\ref{eq:cont-frac-approx}). The interesting constant
$\frac{\pi^2}{6\log 2}$ is sometimes referred to as
the
$\textbf{entropy constant of the continued fraction
map}$.
\end{enumerate}

